% !TeX program = xelatex
% ============================================================
%  第11届信也科技杯 复赛技术报告（≤3 页）
%  ⚠️ 必须用 XeLaTeX 编译；Overleaf 需在 Menu → Compiler 选 XeLaTeX，
%     否则会报 “fandol unavailable / Font unihei58 not found”（那是 pdfLaTeX 模式）。
%  编译方式：XeLaTeX（中文需要 ctex / xeCJK）
%    xelatex finvcup_report.tex
%    xelatex finvcup_report.tex   % 跑两次让引用/目录稳定
%  Overleaf：右上角 Menu → Compiler 选 "XeLaTeX"
% ============================================================
\documentclass[10pt,twocolumn,a4paper]{ctexart}

\usepackage[margin=1.6cm]{geometry}
\usepackage{amsmath,amssymb}
\usepackage{booktabs}
\usepackage{xcolor}
\usepackage{graphicx}
\usepackage{tikz}
\usetikzlibrary{arrows.meta,positioning}
\usepackage[hidelinks]{hyperref}

\setlength{\columnsep}{0.7cm}
\setlength{\parskip}{2pt}

% 紧凑标题
\usepackage{titlesec}
\titlespacing*{\section}{0pt}{6pt}{2pt}
\titlespacing*{\subsection}{0pt}{4pt}{1pt}
\titleformat{\section}{\bfseries\large}{\thesection}{0.5em}{}
\titleformat{\subsection}{\bfseries\normalsize}{\thesubsection}{0.5em}{}

\title{\vspace{-1.2cm}\textbf{多模态融合与集成投票的对话轮次交互建模}\\
\large 第11届信也科技杯 · Turn-Taking 复赛技术报告}
\author{队名：没想好\quad 日期：\today}
\date{}

\begin{document}
\maketitle
\vspace{-0.8cm}

\begin{abstract}
\noindent 本方案面向多方言中文对话的轮次交互（Turn-Taking）建模：给定过去 $\le$30\,s 的双声道
对话（音频、ASR 文本、历史 chunk 标签），预测未来 2\,s（$25\times80$\,ms chunk）内是否出现
延续(C)、静音(NA)、打断(I)、附和(BC)、话权转移(T) 五类事件。我们构建了
\textbf{音频/文本/历史标签/手工统计}四路模态的多模态模型，以\textbf{低秩张量融合 + 自适应门控}
聚合，多标签 sigmoid 输出；训练采用 focal 损失、按类上限的正样本加权、标签平滑与\textbf{权重 EMA}；
推理阶段对 valid 上最优的 \textbf{Top-5} 模型逐标签用各自最优阈值二值化后做\textbf{多数投票}，
以提升在私域测试集（上下文 (0,30]\,s 动态变长）上的鲁棒性。
\end{abstract}

% ---------- 全宽流程图 ----------
\begin{figure*}[t]
\centering
\footnotesize
\begin{tikzpicture}[
  node distance=6mm,
  box/.style={draw,rounded corners,align=center,inner sep=3pt,minimum height=8mm,font=\footnotesize},
  >={Stealth[length=2mm]}]
\node[box,fill=blue!6] (raw) {双声道音频\\ASR 文本\\逐 chunk 标签};
\node[box,fill=blue!6,right=of raw] (slice) {样本切分\\30s 上下文\\+2s 预测窗口\\(动态上下文增强)};
\node[box,fill=green!8,right=of slice] (enc) {四路编码\\Whisper / Qwen3\\标签 CNN / 手工特征};
\node[box,fill=orange!10,right=of enc] (fuse) {低秩张量融合\\+ 自适应门控};
\node[box,fill=orange!10,right=of fuse] (head) {5 路 sigmoid 头\\(多标签)};
\node[box,fill=red!8,right=of head] (ens) {Top-5 成员\\各自阈值二值化\\逐标签多数投票};
\node[box,fill=gray!12,right=of ens] (out) {submit.csv\\segment\_id,\\c,na,i,bc,t};
\draw[->] (raw)--(slice); \draw[->] (slice)--(enc); \draw[->] (enc)--(fuse);
\draw[->] (fuse)--(head); \draw[->] (head)--(ens); \draw[->] (ens)--(out);
\end{tikzpicture}
\caption{从原始数据到最终结果的整体流程：数据切分 $\rightarrow$ 多模态编码 $\rightarrow$ 融合 $\rightarrow$ 多标签预测 $\rightarrow$ 集成投票 $\rightarrow$ 输出。}
\label{fig:pipeline}
\end{figure*}

\section{任务与整体流程}
轮次交互的核心是判断 AI「何时开口、何时收声」。我们将其建模为\textbf{事件级多标签二分类}：
对未来 2\,s 窗口内 5 类标签 $\{C,NA,I,BC,T\}$ 分别预测「是否出现」。整体流程见图~\ref{fig:pipeline}，
满足\textbf{因果约束}——仅使用上下文，不读取未来音频或未来标签。

\section{数据构造}
训练数据为整通双声道对话。以 80\,ms 为一个 chunk，沿对话以步长 \texttt{stride} 滑窗切样本：
取结束位置前 \texttt{context\_chunks}$=375$（30\,s）作为上下文，其后 \texttt{target\_chunks}$=25$（2\,s）
作为预测窗口；标签向量 $y_k=\mathbb{1}[\exists\, \text{未来窗口内出现第 }k\text{ 类}]$。
为贴合私域测试集上下文 $(0,30]$\,s 的\textbf{动态变长}特性，训练时以概率
\texttt{context\_prob} 随机将上下文长度采样到 $[125,375]$ chunk，不足处以 NA(=4) 左侧补齐；
测试时同样把变长上下文归一化到定长 375（截取末段、不足补 NA），保证 batch 内长度统一。

\section{多模态模型}
\subsection{音频编码器}
采用 \textbf{Whisper-large-v3} 的 encoder：整体冻结、仅解冻\textbf{最后 2 层}做任务自适应。
为省显存，冻结前缀在 \texttt{no\_grad} 下前向、仅尾部可训练层建图。在时间轴上只对\textbf{尾部}
（\texttt{tail\_ratio}=0.25）做注意力池化，聚焦临近预测窗口的声学线索，再投影到 512 维。

\subsection{文本编码器}
ASR 文本按说话人加 \texttt{[SPK1]/[SPK2]} 标记拼接，送入冻结的 \textbf{Qwen3-0.6B}；
对序列尾部做带 mask 的注意力池化（近因 utterance 更重要）。tokenizer 采用\textbf{左侧截断/左侧补齐}，
保留最近上下文。

\subsection{历史标签编码与手工特征}
历史 chunk 标签经 Embedding 后走\textbf{双分支 CNN}：尾部分支（最近 \texttt{tail\_k}=75 chunk）
做细粒度卷积+展平，全局分支做卷积+注意力池化，拼接得到标签时序表征。
同时计算\textbf{手工统计特征}：近 25/50/100 窗口的标签分布、到最近事件/话权转移的距离、
标签转移概率、指数/线性时间衰减分布、均值方差变化率、话轮间隔等，投影为 64 维，
为模型注入强先验。

\subsection{低秩张量融合与自适应门控}
四路模态 $\{a,t,c,h\}$ 先做\textbf{低秩张量融合}（rank=64）以紧凑地建模高阶交互：
\[
\textstyle z_{\text{int}}=\Big(\sum_{r=1}^{R}\bigodot_{m}\,\tilde{x}_m W_m^{(r)}\Big)\odot w_f + b,
\]
其中 $\tilde{x}_m=[1;x_m]$。再对每路做投影并由\textbf{门控网络}输出 4 个 $[0,1]$ 权重，
按门控加权后与 $z_{\text{int}}$ 拼接、投影到 320 维融合表征，使模型可按样本自适应地决定信任哪个模态。

\subsection{预测头与损失}
融合表征经线性头输出 5 维 logits，sigmoid 后与标签做 BCE。针对类别不均衡，
采用 \textbf{focal} 调制（$\gamma=2$）+ \textbf{按类上限正样本加权}（cap=8）+ 标签平滑（0.05）：
\[
\mathcal{L}=\frac{1}{K}\sum_{k}\, \alpha_k(1-p_{t,k})^{\gamma}\,\mathrm{BCE}(p_k,\tilde{y}_k).
\]

\section{训练策略}
4$\times$L20 四卡 DDP，等效批 512；峰值 lr $1.2\!\times\!10^{-4}$、warmup 0.03、cosine 衰减、
weight decay 0.03、梯度裁剪 1.0、AMP 混合精度。引入\textbf{权重 EMA}（decay 0.999）：
评估与保存集成成员时使用 EMA 权重，有效闭合大批量带来的泛化间隙。
按 valid \texttt{best\_f1} 选优并早停（patience 15）。

\section{集成与阈值}
训练过程中按 valid \texttt{best\_f1} 排名保留 \textbf{Top-5} 成员；每个成员在 valid 上
\textbf{逐标签}搜索使 F1 最大的阈值，连同成员清单写入 \texttt{ensemble\_manifest.json}。
推理时\textbf{逐成员}（同一时刻 GPU 上仅驻留一个模型，显存峰值=单模型）用各自最优阈值把每个标签
二值化，再做\textbf{逐标签多数投票}：得票 $>$ 半数（5 个成员即 $\ge 3$ 票）判为 1。
成员 checkpoint 为「瘦身」格式（仅存可训练参数），推理时从本地预训练骨干以 \texttt{strict=False}
叠加还原，节省体积。

\section{复现说明}
代码结构与运行命令（详见 \texttt{README.md}）：
\begin{itemize}\setlength{\itemsep}{0pt}
\item \textbf{训练}：\texttt{bash train.sh}（4 卡 DDP，产出 Top-5 成员与 manifest）。
\item \textbf{推理}：\texttt{bash run.sh}（入口），读 \texttt{/xydata} $\rightarrow$ 集成投票 $\rightarrow$
写 \texttt{/app/submit/submit.csv}（列：\texttt{segment\_id,c,na,i,bc,t}）。
\item \textbf{环境}：\texttt{pytorch/pytorch:2.5.1-cuda12.4}，transformers 4.57.x，骨干权重离线本地加载。
\end{itemize}
单成员参数量约 1.2B（Whisper encoder + Qwen3-0.6B + 轻量融合头），推理一次仅载入一个成员，
满足参数量与显存约束。

\end{document}
